\documentclass{article}
\usepackage[utf8]{inputenc}
\usepackage[T1]{fontenc}
\usepackage[french]{babel}
\usepackage{geometry}
\geometry{a4paper, margin=2cm}
\usepackage{hyperref}
\usepackage{listings}
\usepackage{xcolor}

\lstset{
    language=PHP,
    basicstyle=\ttfamily\small,
    keywordstyle=\color{blue},
    commentstyle=\color{green!60!black},
    stringstyle=\color{red},
    breaklines=true,
    captionpos=b,
    numbers=left,
    numberstyle=\tiny,
    frame=single
}

\title{Documentation des Fonctionnalités et Cycle de Vie du Projet de Chiffrement Médical}
\author{Étudiant Ingénieur - Abulcasis, Rabat \\ Stage à la Fondation Cheikh Zayed}
\date{\today}

\begin{document}

\maketitle

\begin{abstract}
Ce document décrit de manière exhaustive les fonctionnalités, concepts et cycle de vie du projet de chiffrement des données médicales pour la Fondation Cheikh Zayed. Il couvre tout le processus de A à Z, incluant les améliorations suggérées pour la mise en production.
\end{abstract}

\tableofcontents
\newpage

\section{Cycle de Vie du Projet : De A à Z}

\subsection{Phase de Conception (A - Analyse et Design)}
Le projet débute par l'identification des besoins en sécurité pour les données médicales sensibles. Les concepts clés incluent :
\begin{itemize}
    \item \textbf{Patron Strategy} : Séparation des algorithmes de chiffrement (AES, RSA) pour une interchangeabilité dynamique.
    \item \textbf{Chiffrement Symétrique (AES-256-GCM)} : Pour rapidité et authenticité (mode AEAD).
    \item \textbf{Chiffrement Asymétrique (RSA-4096)} : Pour distribution sécurisée des clés.
    \item \textbf{Chiffrement Hybride} : Combinaison AES + RSA pour les gros fichiers médicaux.
    \item \textbf{Injection de Dépendances} : Gestion des clés via conteneur Symfony.
    \item \textbf{Transparence ORM} : Chiffrement automatique via Doctrine Event Subscribers.
\end{itemize}

\subsection{Phase de Développement (B - Implémentation)}
Implémentation en Symfony 8 avec PHP 8.4 :
\begin{itemize}
    \item \textbf{EncryptionService} : Service central gérant les stratégies via tagged iterators.
    \item \textbf{Strategies} : AesEncryptionStrategy et RsaEncryptionStrategy.
    \item \textbf{Attribute Encrypted} : Annotation pour marquer les champs sensibles.
    \item \textbf{EncryptionSubscriber} : Hooks Doctrine pour chiffrement/déchiffrement transparent.
    \item \textbf{Entity Patient} : Exemple avec champs chiffrés (socialSecurityNumber en AES, medicalRecord en RSA).
\end{itemize}

\subsection{Phase de Configuration (C - Setup)}
\begin{itemize}
    \item Génération de clés RSA via OpenSSL.
    \item Variables d'environnement (.env) pour clés AES et chemins RSA.
    \item Configuration Doctrine et services Symfony.
\end{itemize}

\subsection{Phase de Test (D - Validation)}
\begin{itemize}
    \item \textbf{TestEncryptionController} : Test unitaire via route Symfony.
    \item \textbf{test\_server.php} : Serveur de démonstration pour chiffrement texte/fichiers.
    \item File d'attente de jobs et worker asynchrone pour le chiffrement/déchiffrement différé de fichiers.
    \item Streaming pour traitement de fichiers lourds sans charger entièrement la mémoire.
    \item Génération de fichiers de test médicaux réalistes (generate\_test\_files.php).
    \item Interface web pour upload/chiffrement/déchiffrement.
\end{itemize}

\subsection{Phase de Déploiement (E - Production)}
\begin{itemize}
    \item \textbf{Dockerfile} : Conteneurisation avec PHP 8.4.
    \item Serveur de test local (php -S).
    \item Préparation pour déploiement cloud (e.g., Render).
\end{itemize}

\subsection{Phase d'Amélioration (F - Optimisations pour Production)}
Améliorations suggérées pour sécurité, flexibilité et conformité :
\begin{itemize}
    \item Correction du padding RSA (PKCS1 vers OAEP).
    \item Gestion avancée des clés (KMS/Vault abstraction, fournisseurs locaux et prêts pour HSM).
    \item Enveloppe de payload chiffrée structurée et versionnée.
    \item Streaming pour gros fichiers et pipeline de traitement sécurisé.
    \item File d'attente de jobs asynchrone et worker de traitement de fichiers.
    \item Authentification et audit logging.
    \item Tests unitaires avec PHPUnit.
    \item Conteneurisation durcie et CI/CD GitHub Actions.
    \item Préparation Kubernetes et déploiement cloud natif.
    \item Conformité RGPD/Maroc.
\end{itemize}

\section{Fonctionnalités Existantes}
\begin{itemize}
    \item Chiffrement/déchiffrement de texte et fichiers.
    \item Support multi-algorithmes (AES/RSA).
    \item Intégration ORM transparente.
    \item Interface de démonstration web.
    \item Génération automatique de clés.
    \item Fichiers de test médicaux.
\end{itemize}

\section{Concepts Clés}
\begin{itemize}
    \item \textbf{Confidentialité} : Données illisibles sans clé.
    \item \textbf{Intégrité/Authenticité} : Via GCM et signatures potentielles.
    \item \textbf{Performance} : AES rapide, RSA pour clés seulement.
    \item \textbf{Évolutivité} : Pattern Strategy extensible.
    \item \textbf{Sécurité} : Défense en profondeur (chiffrement + accès).
\end{itemize}

\section{Archive des Tâches et Améliorations}
Cette section archive les tâches assignées et réalisées pour améliorer le projet.

\subsection{Tâche 1 : Correction du Padding RSA}
\textbf{Description} : Remplacer PKCS1 par OAEP pour renforcer la sécurité RSA. \\
\textbf{Statut} : Implémenté (RSA-OAEP appliqué dans `src/Encryption/Strategy/RsaEncryptionStrategy.php`). \\
\textbf{Date} : \today \\
\textbf{Détails} : `RsaEncryptionStrategy` utilise désormais OAEP padding (`OPENSSL_PKCS1_OAEP_PADDING`), `random_bytes()` pour la génération de clés AES/IV, et effacement mémoire (`sodium_memzero`) lorsque disponible. Voir `src/Encryption/Strategy/RsaEncryptionStrategy.php`.

\subsection{Tâche 2 : Enveloppe AES et abstraction de gestion des clés}
\textbf{Description} : Renforcer la structure des payloads chiffrés et préparer l'intégration KMS/Vault. \\
\textbf{Statut} : Implémenté partiellement (enveloppe versionnée et `LocalKeyProvider` abstrait l'accès aux clés). \\
\textbf{Date} : \today \\
\textbf{Détails} : `AesEncryptionStrategy` et `RsaEncryptionStrategy` utilisent désormais `src/Encryption/Model/EncryptedPayload.php` pour un format payload structuré et versionné. `src/Secrets/LocalKeyProvider.php` implémente `KeyManagementInterface`, et `src/Secrets/VaultKeyProvider.php` expose une architecture prête pour HashiCorp Vault.

\subsection{Tâche 3 : Conteneurisation et pipeline CI/CD}
\textbf{Description} : Mettre en place un Dockerfile durci et un workflow de build/test automatique. \\
\textbf{Statut} : Implémenté (PHP 8.4, extensions de production, non root, GitHub Actions CI, manifests Kubernetes de base). \\
\textbf{Date} : \today \\
\textbf{Détails} : `Dockerfile` a été mis à jour pour une image de production PHP 8.4 avec `mbstring`, `curl`, `zip`, `HEALTHCHECK` et utilisateur non privilégié. `/.github/workflows/ci.yml` contient des étapes de validation composer, tests PHPUnit, analyse de sécurité et vérification de syntaxe. `k8s/` stocke des manifests de déploiement, service et secret pour l'orchestration cloud.

\end{document}